%% ma: L12-B19-0012
%% lop: 12
%% chuong: 6
%% bai: 19
%% dang: 12.19.1
%% loai: DS
%% muc: [NB, TH, VD, VD]
%% dap_an: "ĐĐSS"
%% nguon: "(NBV)-ĐỀ SỐ 7-MỖI NGÀY 1 ĐỀ THI 2025.pdf (D07, MỖI NGÀY 1 ĐỀ - Đề số 7), Phần II Câu 3"
%% kien_thuc: "Bài 19 (công thức xác suất toàn phần), Bài 18 (xác suất có điều kiện)"
%% trang_thai: chinh-thuc
%% ghi_chu: "Trùng ý tưởng với dạng 12.19.1 (cùng mô hình xác suất toàn phần với hai nhóm). Bỏ ảnh minh họa bác sĩ và người cao tuổi (không có dữ kiện toán). Đáp án gốc Đ Đ S S đúng."
\begin{ex}
Khi điều tra sức khỏe nhiều người cao tuổi ở một địa phương, người ta thấy rằng có $40\%$ người cao tuổi bị bệnh tiểu đường. Bên cạnh đó, số người bị bệnh huyết áp cao trong những người bị bệnh tiểu đường là $70\%$, trong những người không bị bệnh tiểu đường là $25\%$. Chọn ngẫu nhiên $1$ người cao tuổi để kiểm tra sức khỏe.
\choiceTF{\True Xác suất chọn được người bị bệnh tiểu đường là $0{,}4$}{\True Xác suất chọn được người bị bệnh huyết áp cao, biết người đó bị bệnh tiểu đường, là $0{,}7$}{Xác suất chọn được người bị bệnh huyết áp cao, biết người đó không bị bệnh tiểu đường, là $0{,}75$}{Xác suất chọn được người bị bệnh huyết áp cao là $0{,}8$}
\loigiai{Gọi $A$: "người bị bệnh tiểu đường", $B$: "người bị bệnh huyết áp cao". Có $P(A)=0{,}4$, $P(B\mid A)=0{,}7$, $P(B\mid\overline A)=0{,}25$.

a) Đúng. b) Đúng. c) Sai: $P(B\mid\overline A)=0{,}25$, không phải $0{,}75$ ($0{,}75$ là xác suất không bị huyết áp cao).

d) Sai: $P(B)=P(A)P(B\mid A)+P(\overline A)P(B\mid\overline A)=0{,}4\cdot0{,}7+0{,}6\cdot0{,}25=0{,}43\ne0{,}8$.}
\end{ex}
